% Generated draft: requires mathematical and source review.
\documentclass[11pt]{article}
\usepackage[margin=25mm]{geometry}
\usepackage{fontspec}
\setmainfont{Latin Modern Roman}
\usepackage{amsmath,amssymb,mathtools,unicode-math}
\setmathfont{Latin Modern Math}
\usepackage{enumitem,needspace,xurl,hyperref}
\setlength{\parindent}{0pt}
\setlength{\parskip}{5pt}
\setlength{\emergencystretch}{3em}
\newcommand{\N}{\mathbb{N}}
\newcommand{\Z}{\mathbb{Z}}
\newcommand{\Q}{\mathbb{Q}}
\newcommand{\R}{\mathbb{R}}
\newcommand{\C}{\mathbb{C}}
\newcommand{\F}{\mathbb{F}}
\newcommand{\A}{\mathbb{A}}
\newcommand{\PP}{\mathbb{P}}
\newcommand{\E}{\mathbb{E}}
\newcommand{\eps}{\varepsilon}
\newcommand{\dd}{\,\mathrm{d}}
\begin{document}
\section*{UnsolvedMath ID: 11100013}
\textbf{Source problem code:} catalogue entry
\subsection{Definitions and mathematical statement}
Fix a prime \(p\) and height \(n\). Let \(K(n)\) be Morava \(K\)-theory and \(L_{K(n)}\) Bousfield localization at it. Choose a precise Morava \(E\)-theory \(E\), its profinite stabilizer group \(S\) and continuous action on \(E\). The source expressly permits either \(L_{K(n)}E(n)\) or a Lubin--Tate theory, and ideally asks for every reasonable choice; these choices must not be identified without comparison. For a Lubin--Tate theory the coefficient topology comes from its complete deformation ring; an extended stabilizer may include Galois automorphisms, depending on the chosen residue field and theory.

Distinguish the stable operation ring \(E^*E=[E,\Sigma^*E]\), ordinary cooperations \(E_*E=\pi_*(E\wedge E)\), and completed cooperations
\[
 E_*^\vee E=\pi_*L_{K(n)}(E\wedge E).
\]
The concrete task is to define these structures and prove the source's expected identifications, with grading/topology appropriate to each choice:
\[
 E^*E\cong E^*[[S]]_{\mathrm{tw}},\qquad
 E_*^\vee E\cong C_{\mathrm{cts}}(S,E_*).
\]
The twisted multiplication is specified on elementary symbols by
\((a[s])(b[t])=a\,s(b)[st]\), reflecting the stabilizer action. Completion affects both coefficients and the profinite group: an inverse-limit description uses compatible coefficient quotients and open finite-index subgroups, or cofinal open normal subgroups with actions descending to the quotients. The function object must use the specified coefficient topology degreewise. Which stabilizer, ideal-adic completion, left/right structures and comparison maps apply are part of the problem, not extra established hypotheses.

More broadly, clarify how these stabilizer structures relate to the \(K(n)\)-local category. The source calls the formulas folklore, credits Hopkins and others with likely knowledge, and discusses Hovey--Strickland's memoir without claiming a published proof there. The summary labels the formulas standard/solved, but its first formula uses cooperation notation for the operation-group-ring assertion. This mismatch is not adopted.
\subsection{Short English statement}
Make precise and prove the stabilizer descriptions of Morava \(E\)-theory operations and completed cooperations, distinguishing the theory choices, coefficient/group completions and continuous functions.
\subsection{Sources}
\begin{itemize}
\item Mark Hovey, \emph{Algebraic Topology Problem List}, ``Morava K- and E-theory,'' item 3, surviving mirror \url{https://www-users.cse.umn.edu/~tlawson/hovey/}. Recorded access 2026-07-29; extraction html-ordered-list.
\item The statement credits Hopkins and others with folklore knowledge and mentions Neil Strickland and Hovey's memoir on the \(K(n)\)-local category as a possible earlier venue, not a supplied proof citation. Exact bibliography entries are absent.
\item Local UnsolvedMath registry, \texttt{data/unsolvedmath/problems.json}, permanent ID \(11100013\), exact source code \texttt{catalogue entry}; supplied solved/SOLVED-IN-LITERATURE summary gives the operation/cooperation descriptions without a precise later reference.
\end{itemize}
\end{document}
